\subsection{INTEGRATION ON LIE ALGEBRAS}

PROPOSITION 2. Let H be a (real) Lie group of dimension m, h its Lie algebra. Let \(\omega_{H}\) be a right-invariant differential form of degree m on H, and let \(\omega_{h}\) be the translation-invariant differential form on h, of degree m, that coincides with \(\omega_{\mathrm{H}}(e)\) at the origin. We have

\[
(\exp_ {\mathrm{H}}) ^ {*} \omega_ {\mathrm{H}} = \lambda_ {\mathfrak {h}} \omega_ {\mathfrak {h}}\tag{10}
\]

where \(\lambda_{\mathfrak{h}}\) is the Ad(H)-invariant function on \(\mathfrak{h}\) such that

\[
\lambda_ {\mathfrak {h}} (x) = \det \left(\sum_ {p \geq 0} \frac {1}{(p + 1) !} (\operatorname{ad} x) ^ {p}\right) \quad \text { for } x \in \mathfrak {h}.
\]

Let \(x, x_1, \ldots, x_m\) be elements of \(\mathfrak{h}\) . We have

\[
(\exp^ {*} \omega_ {\mathrm{H}}) _ {x} (x _ {1}, \dots , x _ {m}) = (\omega_ {\mathrm{H}} (\exp x)) (\mathrm{T} _ {x} (\exp) (x _ {1}), \dots , \mathrm{T} _ {x} (\exp) (x _ {m})).
\]

Denote by \(\varpi(x) : \mathfrak{h} \to \mathfrak{h}\) the right differential of the exponential at \(x\) (Chap. III, §3, no. 17, Def. 8); by definition,

\[
\mathrm{T} _ {x} (\exp) (y). (\exp x) ^ {- 1} = \varpi (x). y \quad \text { for   all } y \in \mathfrak {h}.
\]

The form \(\omega_{H}\) being right invariant, we obtain

\[
\begin{array}{r l} & (\omega_ {\mathrm{H}} (\exp x)) (\mathrm{T} _ {x} (\exp) (x _ {1}), \ldots , \mathrm{T} _ {x} (\exp) (x _ {m})) \\ & \quad = \omega_ {\mathrm{H}} (e) (\varpi (x). x _ {1}, \ldots , \varpi (x). x _ {m}) = (\det \varpi (x)) \omega_ {\mathfrak {h}} (x _ {1}, \ldots , x _ {m}); \end{array}
\]

thus, \(\exp^{*}\omega_{\mathrm{H}} = \lambda_{\mathfrak{h}}\omega_{\mathfrak{h}}\) , with \(\lambda_{\mathfrak{h}}(x) = \det \varpi(x) = \det \frac{\exp\operatorname{ad} x - 1}{\operatorname{ad} x}\) (Chap. III, §6, no. 4, Prop. 12).

Let \(h \in H\) ; since Ad h is an automorphism of h, we have

\[
\operatorname{ad} \left((\operatorname{Ad} h) (x)\right) = \operatorname{Ad} h \circ \operatorname{Ad} x \circ (\operatorname{Ad} h) ^ {- 1},
\]

so \(\lambda_{\mathfrak{h}}((\mathrm{Ad}h)(x)) = \lambda_{\mathfrak{h}}(x)\) . Thus, the function \(\lambda_{\mathfrak{h}}\) is invariant under \(\mathrm{Ad(H)}\) ; this completes the proof of the proposition.

Remark. Consider the function \(\lambda_{g}\) associated to a compact Lie group G; in view of §2, no. 1, Th. 1, to calculate \(\lambda_{g}\) it suffices to know its values on t. But, for \(x \in t\) , the endomorphism ad x of g is semi-simple, and has eigenvalues 0 (with multiplicity r) and, for all \(\alpha \in \mathrm{R}(\mathrm{G}, \mathrm{T})\) , \(\delta(\alpha)(x)\) (with multiplicity 1). It follows immediately that

\[
\lambda_ {\mathfrak {g}} (x) = \prod_ {\alpha \in \mathrm{R} (\mathrm{G}, \mathrm{T})} \frac {e ^ {\delta (\alpha) (x)} - 1}{\delta (\alpha) (x)} = \frac {\delta_ {\mathfrak {g}} (x)}{\pi_ {\mathfrak {g}} (x)}\tag{11}
\]

\[
\text { with } \delta_ {\mathfrak {g}} (x) = \delta_ {\mathrm{G}} (\exp x) \text { and } \pi_ {\mathfrak {g}} (x) = \prod_ {\alpha \in \mathrm{R} (\mathrm{G}, \mathrm{T})} \delta (\alpha) (x) = \det \operatorname{ad} _ {\mathfrak {g} / \mathfrak {t}} (x).
\]

Let \(\omega_{\mathrm{G / T}}\) be an invariant differential form of degree \(n - r\) on \(\mathrm{G / T}\) and \(\omega_{\mathfrak{t}}\) a translation-invariant differential form of degree \(r\) on \(\mathfrak{t}\) . With the notation of no. 1, denote by \(\omega_{\mathrm{G / T}} \cap \omega_{\mathfrak{t}}\) the unique translation-invariant differential form \(\omega_{\mathfrak{g}}\) of degree \(n\) on \(\mathfrak{g}\) such that \(\omega_{\mathfrak{g}}(0) = \omega_{\mathrm{G / T}}(\bar{e}) \cap \omega_{\mathfrak{t}}(0)\) .

Finally, denote by \(\psi:(\mathrm{G}/\mathrm{T})\times\mathfrak{t}\to\mathfrak{g}\) the morphism of manifolds induced by the map \((g,x)\mapsto(\operatorname{Ad}g)(x)\) from \(G\times t\) to g by passage to the quotient.

PROPOSITION 3. Let \(\omega_{\mathfrak{g}}\) , \(\omega_{\mathfrak{t}}\) , \(\omega_{\mathrm{G / T}}\) be invariant differential forms on \(\mathfrak{g}\) , \(\mathfrak{t}\) , \(\mathrm{G / T}\) , respectively, of respective degrees \(n, r, n - r\) . If \(\omega_{\mathfrak{g}} = \omega_{\mathrm{G / T}} \cap \omega_{\mathfrak{t}}\) , we have

\[
\psi^ {*} \omega_ {\mathfrak {g}} = \omega_ {\mathrm{G/T}} \wedge \pi_ {\mathfrak {g}} \omega_ {\mathfrak {t}}
\]

where \(\pi_{\mathfrak{g}}\) is the function on \(\mathfrak{t}\) defined by \(\pi_{\mathfrak{g}}(x) = \prod_{\alpha \in \mathrm{R}(\mathrm{G},\mathrm{T})}\delta (\alpha)(x)\) .

Denote by \(\omega_{G}\) (resp. \(\omega_{T}\) ) the invariant differential form of maximum degree on G (resp. T) that coincides with \(\omega_{g}\) (resp. \(\omega_{t}\) ) at the origin. Consider the commutative diagram

\[
\begin{array}{c c c} (\mathrm{G} / \mathrm{T}) \times \mathfrak {t} & \stackrel {{\psi}} {{\longrightarrow}} & \mathfrak {g} \\ \Big \downarrow^ {(\mathrm{Id}, \exp_ {\mathrm{T}})} & & \Big \downarrow^ {\exp_ {\mathrm{G}}} . \\ (\mathrm{G} / \mathrm{T}) \times \mathrm{T} & \stackrel {{f}} {{\longrightarrow}} & \mathrm{G} \end{array}
\]

In view of Prop. 1 of no. 2 and the relation \(\exp_{T}^{*}\omega_{T} = \omega_{t}\) , we deduce the equality

\[
\psi^ {*} \mathrm{exp} _ {\mathrm{G}} ^ {*} \omega_ {\mathrm{G}} = \omega_ {\mathrm{G/T}} \wedge \delta_ {\mathfrak {g}} \omega_ {\mathfrak {t}}.
\]

By Prop. 2, \(\psi^{*}\exp_{\mathrm{G}}^{*}\omega_{\mathrm{G}} = (\psi^{*}\lambda_{\mathfrak{g}})\psi^{*}\omega_{\mathfrak{g}}\) . Since the function \(\lambda_{\mathfrak{g}}\) is invariant under \(\operatorname{Ad}(\mathrm{G})\) , we have

\[
(\psi^ {*} \lambda_ {\mathfrak {g}}) (\bar {g}, x) = \lambda_ {\mathfrak {g}} (x) = \frac {\delta_ {\mathfrak {g}} (x)}{\pi_ {\mathfrak {g}} (x)} \quad \text { for } \bar {g} \in \mathrm{G/T}, x \in \mathfrak {t}.
\]

It follows that the forms \(\psi^{*}\omega_{\mathrm{G}}(\bar{g},x)\) and \(\omega_{\mathrm{G / T}}(\bar{g})\wedge \pi_{\mathfrak{g}}(x)\omega_{\mathfrak{t}}(x)\) coincide where \(\delta_{\mathfrak{g}}(x)\) is non-zero, that is on the dense open subset \((\mathrm{G / T})\times \mathfrak{t}_r\) ; thus, they are equal, hence the proposition.

Choose invariant differential forms \(\omega_{G}\) on G and \(\omega_{T}\) on T, of maximum degree, such that \(|\omega_{G}| = dg\) and \(|\omega_{T}| = dt\) ; denote by \(\omega_{g}\) (resp. \(\omega_{t}\) ) the translation-invariant differential form on g (resp. t) that coincides with \(\omega_{\mathrm{G}}(e)\) (resp. \(\omega_{\mathrm{T}}(e)\) ) at the origin, and dz (resp. dx) the Haar measure \(|\omega_{g}|\) (resp. \(|\omega_{t}|\) ). Reasoning as in no. 2, mutatis mutandis, gives the following proposition:

PROPOSITION 4. The measure \(dz\) on \(\mathfrak{g}\) is the image under the proper map \((g,x)\mapsto (\operatorname {Ad}g)(x)\) from \(\mathrm{G}\times \mathfrak{t}\) to \(\mathfrak{g}\) of the measure \(dg\otimes \frac{1}{w(\mathrm{G})}\pi_{\mathfrak{g}}dx\) .

We leave to the reader the statement and proof of the analogues of Cor. 1 to 3 and Remarks 1 to 3 of no. 2. For example, let \(\varphi\) be an integrable function on \(\mathfrak{g}\) (with values in a Banach space or \(\overline{\mathbf{R}}\) ); then

\[
\int_ {\mathfrak {g}} \varphi (z) d z = \frac {1}{w (\mathrm{G})} \int_ {\mathfrak {t}} \left(\int_ {\mathrm{G}} \varphi ((\mathrm{Ad} g) x) d g\right) \pi_ {\mathfrak {g}} (x) d x,\tag{12}
\]

and, in particular, if \(\varphi\) is invariant under \(\operatorname{Ad}(\mathbf{G})\) ,

\[
\int_ {\mathfrak {g}} \varphi (z) d z = \frac {1}{w (\mathrm{G})} \int_ {\mathfrak {t}} \varphi (x) \pi_ {\mathfrak {g}} (x) d x.\tag{13}
\]

